\documentclass[11pt,a4paper]{article}
\usepackage[top=1.7cm,bottom=1.5cm,left=1.5cm,right=1.5cm,headsep=8pt,footskip=16pt]{geometry}
\usepackage{enumitem}
\usepackage{amssymb}
\usepackage{array}
\usepackage[table]{xcolor}
\usepackage{fancyhdr}
\usepackage{lastpage}
\usepackage{multicol}
\usepackage{polyglossia}
\setdefaultlanguage[calendar=gregorian, numerals=maghrib]{arabic}
\setotherlanguage{english}
\IfFontExistsTF{Amiri}{%
  \newfontfamily\arabicfont[Script=Arabic, Scale=1.08]{Amiri}
}{%
  \newfontfamily\arabicfont[Script=Arabic, Scale=1.08, Path=../fonts/, Extension=.ttf,
    UprightFont=*-Regular, BoldFont=*-Bold, ItalicFont=*-Italic,
    BoldItalicFont=*-BoldItalic]{Amiri}
}
\IfFontExistsTF{AlBattar}{%
  \newfontfamily\khatt[Script=Arabic, Scale=1.15]{AlBattar}
}{%
  \newfontfamily\khatt[Script=Arabic, Scale=1.15, Path=../fonts/,
    BoldFont=ae_AlBattar.ttf]{ae_AlBattar.ttf}
}
\renewcommand{\headrulewidth}{0.8pt}
\renewcommand{\footrulewidth}{0.8pt}
\fancyhf{}
\rhead{\khatt المدرسة العليا للأساتذة ببوزريعة}
\lhead{\khatt ملحقة بومرداس}
\lfoot{\khatt الأستاذ: د. بن عبد الرزاق يوسف}
\cfoot{\thepage/\pageref{LastPage}}
\pagestyle{fancy}
\parindent 0pt
\setlength{\columnsep}{0.6cm}
\setlength{\columnseprule}{0.4pt}
\setlist[enumerate,1]{label=\textbf{\arabic*.}, leftmargin=*, itemsep=5pt, topsep=2pt}
\begin{document}

\linespread{1.08}\selectfont
\begin{center}
{\LARGE\khatt\textbf{الاختبار الثاني والعشرون في الذكاء الاصطناعي}}\\[4pt]
{\normalsize وضعيات إشكالية ــ أسئلة متعددة الاختيارات ــ الفصول من 1 إلى 6}
\end{center}
\vspace{-2pt}\rule{\textwidth}{0.6pt}\\[2pt]
\textbf{التعليمة:} يتضمن الاختبار خمس وضعيات إشكالية. اقرأ كل وضعية بعناية، ثم اختر التحليل أو التصرف الأنسب
من بين الاختيارات الخمسة؛ لكل وضعية إجابة صحيحة واحدة فقط.
اكتب \textbf{حرف} الإجابة الصحيحة في جدول الإجابات في آخر الاختبار.
لكل وضعية 4 نقاط (المجموع 20 نقطة).
\vspace{4pt}

\noindent
{\large\khatt\textbf{الوضعية 1}}\hfill{\small(الفصل 1)}\par\nopagebreak\vspace{3pt}
\setlength{\fboxsep}{7pt}\colorbox{gray!10}{\parbox{\dimexpr\linewidth-14pt}{يُعدّ نادي المسرح في متوسطة مسرحية قصيرة بالدارجة الجزائرية، فطلب الأستاذ المشرف من أداة توليدية كتابة الحوار. جاء الحوار بلهجة أقرب إلى اللهجات المشرقية، مع كلمات جزائرية قليلة مستعملة في غير موضعها. ويرى أحد التلاميذ أن الأداة «لا تعرف الجزائر». ما التقدير الأدق والتصرف الأنسب؟}}\par\nopagebreak
\begin{itemize}[leftmargin=1.0cm, labelsep=6pt, itemsep=3pt, topsep=4pt, parsep=0pt]
  \item[\textbf{أ)}] الأداة تعرف كل اللهجات العربية بالدقة نفسها، والخلل في صياغة الطلب وحدها، فيكفي أن يكتب «بالدارجة الجزائرية الأصيلة» ليأتي الحوار مطابقًا تمامًا للهجة التلاميذ دون أي تعديل.
  \item[\textbf{ب)}] اللهجات الأقل حضورًا في بيانات التدريب تُنتَج بجودة أضعف، فيستعين بالأداة في بناء المشاهد والأفكار، ثم يكتب التلاميذ الحوار بلهجتهم أو يعدّلونه، لأنهم أعرف بها من الأداة.
  \item[\textbf{ج)}] يكتب المسرحية بالفصحى بدل الدارجة، لأن الأدوات التوليدية لا تستطيع إنتاج أي نص بلهجة عربية، وكل ما تكتبه بغير الفصحى مختلق، فلا يمكن الاستفادة منه حتى في بناء المشاهد.
  \item[\textbf{د)}] يعتمد الحوار كما هو، لأن الجمهور يفهم اللهجات المشرقية من خلال المسلسلات، واستعمال كلمات جزائرية في غير موضعها يضيف طرافة إلى المسرحية، فلا حاجة إلى إشراك التلاميذ في التعديل.
  \item[\textbf{هـ)}] يطلب من الأداة ترجمة الحوار إلى الفرنسية ثم إعادته إلى الدارجة، لأن المرور بلغة وسيطة يجعل الأداة تختار المفردات الجزائرية الصحيحة تلقائيًا، فيأتي الحوار سليمًا دون مراجعة.
\end{itemize}
\par\vspace{6pt}
\noindent
{\large\khatt\textbf{الوضعية 2}}\hfill{\small(الفصل 2)}\par\nopagebreak\vspace{3pt}
\setlength{\fboxsep}{7pt}\colorbox{gray!10}{\parbox{\dimexpr\linewidth-14pt}{يخطط فريق لبناء نموذج يولّد تمارين نحوية، ولتوفير الجهد اقترح عضو تدريبه على عشرات الآلاف من التمارين التي ولّدتها أدوات أخرى دون مراجعة، بدل تمارين كتبها أساتذة. ويرى أن «الكمية ستعوّض كل شيء». ويريد الفريق قرارًا قبل بدء التدريب. ويضيف أن تمارين الأساتذة قليلة ومكلفة. ما التقدير الأدق؟}}\par\nopagebreak
\begin{itemize}[leftmargin=1.0cm, labelsep=6pt, itemsep=3pt, topsep=4pt, parsep=0pt]
  \item[\textbf{أ)}] الاقتراح مقبول بشرط مضاعفة عدد التمارين المولَّدة، لأن الأخطاء تتلاشى مع زيادة الكمية، فلا حاجة إلى المراجعة أو إلى أي تمارين كتبها أساتذة متخصصون في تعليم النحو.
  \item[\textbf{ب)}] الاقتراح جيد إذا خُلطت التمارين المولَّدة بعدد قليل جدًا من تمارين الأساتذة، لأن التمارين البشرية القليلة تصحّح كل أخطاء التمارين المولَّدة الكثيرة دون حاجة إلى مراجعتها.
  \item[\textbf{ج)}] التدريب على مخرجات آلية غير مراجعة ينقل أخطاءها ويضيّق تنوعها، وقد تتراكم الأخطاء جيلًا بعد جيل؛ فالأنسب بيانات بشرية مراجعة، أو مخرجات آلية خضعت لتدقيق صارم.
  \item[\textbf{د)}] الاقتراح ممتاز، لأن مخرجات الأدوات الأخرى خالية من الأخطاء بالضرورة، فهي ناتجة عن نماذج مدربة جيدًا، والكمية الكبيرة تجعل النموذج الجديد أفضل من الأدوات التي ولّدت تمارينه.
  \item[\textbf{هـ)}] الاقتراح مرفوض لأن النماذج لا يمكنها أبدًا أن تتعلم من نصوص مولَّدة، فهي تتجاهلها تلقائيًا أثناء التدريب، وبذلك لا يتعلم النموذج شيئًا مهما كان عدد التمارين المستعملة.
\end{itemize}
\par\vspace{6pt}
\noindent
{\large\khatt\textbf{الوضعية 3}}\hfill{\small(الفصل 3)}\par\nopagebreak\vspace{3pt}
\setlength{\fboxsep}{7pt}\colorbox{gray!10}{\parbox{\dimexpr\linewidth-14pt}{طلب طالب من الأداة كتابة مقال من ألف كلمة عن «الالتزام في الشعر العربي الحديث»، فجاء المقال متماسكًا لغويًا، لكنه مكرّر الأفكار مضطرب الخطة. وأستاذه يشترط أن يكون المقال من إنجازه بمساعدة الأداة لا من إنتاجها، مع الإفصاح عن دورها. ما الطريقة الأنسب للاستعانة بها؟}}\par\nopagebreak
\begin{itemize}[leftmargin=1.0cm, labelsep=6pt, itemsep=3pt, topsep=4pt, parsep=0pt]
  \item[\textbf{أ)}] يطلب من الأداة خطة مفصلة ثم مقالًا كاملًا وفقها، ويسلّمه باسمه، لأن الخطة الواضحة تجعل المقال من إنجازه، والأداة مجرد منفذ لخطة طلبها هو منها بصياغته الخاصة.
  \item[\textbf{ب)}] يكتب المقال بنفسه دون أي استعانة، ثم يطلب من الأداة إعادة كتابته كاملًا بأسلوب أفضل، ويسلّم النسخة المعاد كتابتها، لأن الأفكار الأصلية له مهما تغيّرت الصياغة.
  \item[\textbf{ج)}] يطلب من الأداة خطة ثم يكتب المقال بنفسه، لكنه لا يفصح عن أي استعانة، لأن الخطة لا تُعدّ جزءًا من المقال، والإفصاح مطلوب فقط عندما تكتب الأداة فقرات كاملة، فيبقى عمله أصيلًا في نظر أستاذه.
  \item[\textbf{د)}] يضع خطته وأفكاره بنفسه أولًا، ثم يطلب من الأداة نقدها واقتراح ما ينقصها، ويكتب الفقرات بنفسه، ويستعين بها في مراجعة الأسلوب، مع الإفصاح عن دورها، فيبقى البناء والتفكير من عمله.
  \item[\textbf{هـ)}] يطلب من الأداة كتابة المقال عشر مرات، ويختار من كل نسخة أحسن فقرة، ثم يجمعها في مقال واحد يسلّمه، لأن الاختيار والجمع عمل شخصي يكفي لجعل المقال من إنجازه، ويمثّل جهده الحقيقي في إنجاز العمل.
\end{itemize}
\par\vspace{6pt}
\noindent
{\large\khatt\textbf{الوضعية 4}}\hfill{\small(الفصل 4)}\par\nopagebreak\vspace{3pt}
\setlength{\fboxsep}{7pt}\colorbox{gray!10}{\parbox{\dimexpr\linewidth-14pt}{تداول تلاميذ مقطع فيديو مولَّدًا يظهر فيه أستاذ الرياضيات وهو يرقص ويقول كلامًا لم يقله، صنعه تلميذ «للمزاح» بأداة مجانية. انتشر المقطع خارج المؤسسة، وتأثر الأستاذ نفسيًا. ويطلب المدير رأي الأساتذة في المعالجة قبل اجتماع مجلس التأديب. ويريد بعض الأساتذة استثمار الحادثة تربويًا. ما المعالجة الأنسب؟}}\par\nopagebreak
\begin{itemize}[leftmargin=1.0cm, labelsep=6pt, itemsep=3pt, topsep=4pt, parsep=0pt]
  \item[\textbf{أ)}] يُتعامل مع الفعل بوصفه مساسًا بكرامة شخص وانتحالًا لصورته وفق النظام الداخلي، مع حذف المقطع حيث أمكن، ويُستثمر الحادث في توعية التلاميذ بخطورة التزييف العميق وحقوق الصورة.
  \item[\textbf{ب)}] يُكتفى بحذف المقطع من حسابات التلاميذ دون أي إجراء أو توعية، لأن الفعل وقع خارج القسم، والمؤسسة غير معنية بما ينشره التلاميذ على مواقع التواصل خارج أوقات الدراسة.
  \item[\textbf{ج)}] لا يُعاقب التلميذ، لأن المقطع مزاح واضح لا يصدّقه أحد، والأداة المجانية هي المسؤولة عن إنتاجه، فيكفي أن يعتذر التلميذ للأستاذ شفهيًا دون حذف المقطع أو أي إجراء آخر.
  \item[\textbf{د)}] يُعاقب التلميذ أشد العقوبات دون حوار أو توعية، ويُمنع كل التلاميذ من استعمال الهواتف نهائيًا، لأن التوعية لا تجدي مع هذه الأفعال، والعقاب الجماعي وحده يمنع تكرارها.
  \item[\textbf{هـ)}] يُطلب من الأستاذ أن ينشر مقطعًا مولَّدًا مماثلًا للتلميذ، حتى يفهم التلميذ أثر فعله بنفسه، لأن المعاملة بالمثل أنجع وسيلة تربوية لمعالجة سلوكيات التزييف العميق.
\end{itemize}
\par\vspace{6pt}
\noindent
{\large\khatt\textbf{الوضعية 5}}\hfill{\small(الفصل 6)}\par\nopagebreak\vspace{3pt}
\setlength{\fboxsep}{7pt}\colorbox{gray!10}{\parbox{\dimexpr\linewidth-14pt}{طلب أستاذ من أداة توليد صور خريطة للأندلس في عهد عبد الرحمن الناصر لدرس عن الأدب الأندلسي. جاءت الخريطة جميلة الألوان، لكن أسماء بعض المدن مكتوبة بحروف مشوهة، وحدود الدولة لا تطابق ما في كتب التاريخ، وظهرت مدن في غير مواضعها. ما التصرف الأنسب؟}}\par\nopagebreak
\begin{itemize}[leftmargin=1.0cm, labelsep=6pt, itemsep=3pt, topsep=4pt, parsep=0pt]
  \item[\textbf{أ)}] يطلب من الأداة تصحيح أسماء المدن وحدها، ثم يعتمد الخريطة، لأن الحدود والمواقع لا تتغير في الخرائط المولَّدة، فإذا صحت الأسماء صحت الخريطة كلها دون حاجة إلى مقارنتها بالأطلس، فتصبح الخريطة صالحة للعرض.
  \item[\textbf{ب)}] يكرر الطلب حتى يحصل على خريطة تبدو دقيقة، ثم يعتمدها دون مقارنتها بأي أطلس، لأن الخريطة التي تبدو دقيقة للعين تكون مطابقة للحدود التاريخية بالضرورة، فيكسب وقتًا في التحضير.
  \item[\textbf{ج)}] يعتمد الخريطة كما هي، لأن هدف الدرس أدبي لا جغرافي، والتلاميذ لا يحتاجون إلى دقة الحدود والمواقع، والجمال البصري للخريطة كافٍ لتقريب الأندلس إلى أذهانهم.
  \item[\textbf{د)}] يتخلى عن استعمال أي خريطة في الدرس، لأن كل الخرائط التاريخية، المولَّدة وغير المولَّدة، تقريبية بالضرورة، فلا فرق في الدقة بين أطلس تاريخي محقق وخريطة مولَّدة.
  \item[\textbf{هـ)}] لا يعتمد الخريطة المولَّدة مرجعًا جغرافيًا، بل يستعمل خريطة من أطلس تاريخي موثوق، ويمكنه الاستفادة من الأداة في عناصر تزيينية فقط، لأن أدوات توليد الصور لا تضمن دقة المواقع والأسماء.
\end{itemize}
\par\vspace{6pt}
\begin{center}
{\khatt\textbf{جدول الإجابات}}\\[3pt]
\renewcommand{\arraystretch}{1.6}
\begin{tabular}{|>{\centering\arraybackslash}m{1.9cm}|*{5}{>{\centering\arraybackslash}m{1.2cm}|}}
\hline
\textbf{رقم السؤال} & \textbf{1} & \textbf{2} & \textbf{3} & \textbf{4} & \textbf{5} \\ \hline
\textbf{الإجابة} &  &  &  &  &  \\[8pt] \hline
\end{tabular}

\vspace{8pt}
\renewcommand{\arraystretch}{2.2}
\begin{tabular}{|>{\centering\arraybackslash}m{7.5cm}|>{\centering\arraybackslash}m{3.2cm}|>{\centering\arraybackslash}m{3.2cm}|}
\hline
\textbf{الاسم واللقب} & \textbf{الفوج} & \textbf{العلامة} \\ \hline
 & & \hfill /20 \\ \hline
\end{tabular}
\end{center}
\end{document}
